% ============================================================
% Capitulo 1 -- Introducao
% Tamanho-alvo: 4-6 paginas | Sem figuras nem tabelas
% Ordem de escrita: ULTIMO (depois de saber qual historia contar)
% ============================================================

\chapter{Introdu\c{c}\~ao}
\label{cap:introducao}

% ============================================================
\section{Contextualiza\c{c}\~ao}
\label{sec:contexto}

% TESE: desinformacao e' problema epistemico que antecede LLMs; NLP isolada nao basta.
%
% CONTEUDO (~1 pagina):
% - Por que fake news e' problema (impacto democratico/social).
% - Limites do NLP puro: LLMs geram texto gramaticalmente impecavel.
% - Topologia da propagacao como sinal complementar (nao substituto).
% - Cuidado: nao culpar LLMs pela crise; eles aceleram, nao causam.
%
% CITACOES: \cite{ieee_fakenewslandscape}, \cite{pmc_overviewfakenews},
%           \cite{scielo_posttruth}, \cite{arxiv_llmfakenews}.

% ============================================================
\section{Motiva\c{c}\~ao}
\label{sec:motivacao}

% TESE: topologia e' complementar e independente do idioma do post.
%
% CONTEUDO (~0.5 pagina):
% - Topologia nao processa texto -> mecanismo arquiteturalmente
%   independente de idioma.
% - Aplicacao pratica: cenarios onde texto nao e' confiavel ou nao existe.
%
% CITACAO: \cite{monti2019fakenews} (fundador da linha GNN para fake news).

% ============================================================
\section{Perguntas de pesquisa}
\label{sec:rqs}

% TESE: 3 RQs operacionalizaveis com experimentos diretos.
%
% CONTEUDO (~0.5 pagina) -- listar como perguntas numeradas:
% - RQ1: topologia carrega sinal acima do BERT da raiz?
% - RQ2: sob quais condicoes estruturais? (operacionalizado via Cohen's d)
% - RQ3: e' viavel detectar so com topologia, em cenarios sem texto util?
%        (parte a: experimento estrutural puro; parte b: invariancia a idioma)

% ============================================================
\section{Contribui\c{c}\~oes}
\label{sec:contribuicoes}

% TESE: 4 contribuicoes distintas, todas com artefatos rastreaveis.
%
% CONTEUDO (~1 pagina) -- lista numerada, 1-2 frases cada:
% (a) Ablation sistematica de encodings posicionais (corrige Erro 1).
% (b) Demonstracao positivo+negativo (GossipCop vs PolitiFact) com
%     explicacao estrutural via Cohen's d.
% (c) Analise da concordancia textual x topologico + regra dual post-hoc.
% (d) Pesos no Hugging Face Hub + ferramenta web (FastAPI + Next.js).
%
% FECHAMENTO: citar Apendice~\ref{ap:repo} pra detalhes tecnicos.

% ============================================================
\section{Estrutura do trabalho}
\label{sec:estrutura}

% TESE: mapa dos 5 capitulos restantes.
%
% CONTEUDO (~0.5 pagina) -- 1 paragrafo por capitulo, 2 frases cada:
% Cap 2 (Fundamentacao) -> Cap 3 (Metodologia) ->
% Cap 4 (Resultados) -> Cap 5 (Aplicacao) -> Cap 6 (Conclusao).
